\section{Introduction}

This paper is the third in a sequence studying related finite covers
and quotient structures over the line graph of the Petersen graph.

The first companion paper studies
\[
X=KC(L(P)),
\]
the canonical bipartite double cover of the Petersen line graph. It
shows that the symmetry visible from one selected quotient frame,
\(S_5\times C_2\), is properly contained in the full automorphism
group
\[
\operatorname{Aut}(X)\cong S_5\times S_3.
\]
The complementary \(S_3\) permutes a three-member orbit of regular
two-fold quotient frames, each having quotient isomorphic to
\(L(P)\) \cite{CaveCanonicalBipartite2026}.

The second companion paper determines the full automorphism group of
the distinct sixty-vertex graph
\[
G_{60}\cong \operatorname{AT4val}[60,6].
\]
Using the canonical two-fold quotient \(G_{60}\to G_{30}\), an
independent exhaustive native enumeration, internal group invariants,
and an explicit \(480\)-element multiplication certificate, it proves
\[
\operatorname{Aut}(G_{60})
\cong S_5\times_{C_2}D_8
\]
for the fiber product in which the sign character of \(S_5\) is
matched with a quotient character of \(D_8\) having Klein-four kernel
\cite{CaveAT4Automorphism2026}.

That classification determines which abstract group acts on
\(G_{60}\). The present paper asks why the group has this architecture
in the native graph.

We reconstruct the Klein-four kernel
\[
V_4=\{1,a,b,ab\}
\]
as a free normal group of graph automorphisms. Its orbits give fifteen
four-vertex chambers and the quotient
\[
G_{60}/V_4\cong L(P).
\]
The associated character decomposition separates one distinguished
nontrivial sector \(a\) from an exchanged pair \(b,ab\).

The fifteen chambers further split into five invariant triples. Under
the explicit Petersen coordinate bridge, each triple consists of
three pairwise disjoint Petersen edges. The five triples therefore
form five three-edge matchings partitioning all fifteen Petersen
edges. Every automorphism preserves this block system, and its induced
action is the full symmetric group \(S_5\). The native \(V_4\) is the
exact kernel.

This gives the split exact sequence
\[
1\longrightarrow V_4
\longrightarrow \operatorname{Aut}(G_{60})
\longrightarrow S_5
\longrightarrow 1
\]
and the equivalent semidirect-product description
\[
\operatorname{Aut}(G_{60})
\cong V_4\rtimes_{\operatorname{sgn}}S_5.
\]
Every quotient permutation fixes \(a\). Even permutations fix \(b\)
and \(ab\) separately, while odd permutations exchange them. Thus the
parity coupling appearing abstractly in the earlier fiber-product
classification is realized concretely as conjugation on the native
Klein register.

A second part of the paper studies cycle-level consequences of this
register. Quotient pentagons with nontrivial \(b\)- or \(ab\)-holonomy
lift to paired ten-cycles requiring two quotient traversals for
registered closure. Exactly twenty-four simple ten-cycles have this
form. They constitute one automorphism orbit and admit an exact
fourth-order adjacency and Laplacian signature.

Within the adjacent \(120\)-cycle frontier, the registered cycles are
distinguished by the absence of sixteen external four-step returns.
The frontier cycles carry two disjoint relay paths and have
non-native Klein-four stabilizers. There are fifteen such foreign
Klein groups. They form one conjugacy class, are pairwise disjoint
away from the identity, and correspond to the three internal pairs
within each of the five Petersen matching blocks.

The resulting structure provides an equivariant interpretation of
\[
60=5\cdot3\cdot4.
\]
The factor \(5\) records the invariant Petersen matching blocks, the
factor \(3\) records the chambers within each block, and the factor
\(4\) records the native Klein states within each chamber.

The paper is organized as follows. Section 2 records the precise
dependency on the two companion classifications. Sections 3 and 4
recover the native Klein register, quotient, and character operators.
Sections 5 through 7 develop cycle holonomy, registered closure, and
the fourth-moment frontier. Sections 8 through 10 classify the foreign
Klein stabilizers, Petersen matchings, and canonical five-point
register. Section 11 proves the split extension and parity action.
Sections 12 and 13 collect the principal theorem and its finite
certificates, and Section 14 concludes.
